\chapter{Lifting of Quantum Torsion}
\label{chap:torsion}

% ── Introduction ──────────────────────────────────────────────────────────────
\section{Introduction}
\label{sec:torsion-intro}

Torsion is a classical differential-geometric concept measuring the failure of
a connection to be symmetric.  In this chapter we introduce a quantum analogue,
\emph{quantum torsion}, and study its \emph{lifting} to higher-categorical and
operator-algebraic structures within the QOU framework developed in
Chapter~\ref{chap:qou}.

% ── Quantum Torsion ───────────────────────────────────────────────────────────
\section{Quantum Torsion}
\label{sec:quantum-torsion}

\begin{definition}[Quantum Connection]
  A \emph{quantum connection} on a Hilbert bundle
  $\mathcal{E} \to \mathcal{M}$ is a $\mathbb{C}$-linear map
  $\nabla: \Gamma(\mathcal{E}) \to \Gamma(T^*\mathcal{M} \otimes \mathcal{E})$
  satisfying the Leibniz rule with respect to the $C^\infty(\mathcal{M})$
  module structure.
\end{definition}

\begin{definition}[Quantum Torsion Tensor]
  The \emph{quantum torsion tensor} $\mathcal{T}^\nabla$ associated to a
  quantum connection $\nabla$ is defined by
  \begin{equation}
    \mathcal{T}^\nabla(X, Y)
    = \nabla_X Y - \nabla_Y X - [X, Y],
    \quad X, Y \in \Gamma(T\mathcal{M}).
    \label{eq:torsion}
  \end{equation}
\end{definition}

% ── Lifting ───────────────────────────────────────────────────────────────────
\section{Lifting to Higher Structures}
\label{sec:lifting}

% TODO: define the lifting functor and prove its existence.

\begin{definition}[Torsion Lifting]
  Let $\mathcal{C}$ be a monoidal category of quantum field configurations.
  A \emph{torsion lifting} is a functor
  $\mathcal{L}: \mathcal{C} \to \mathcal{C}^{(2)}$
  to the two-categorical extension $\mathcal{C}^{(2)}$ such that the torsion
  tensor $\mathcal{T}^\nabla$ lifts coherently along $\mathcal{L}$.
\end{definition}

\begin{theorem}[Existence of Torsion Lifting]
  \label{thm:lifting-exists}
  Under suitable curvature conditions on $\nabla$, the torsion lifting
  $\mathcal{L}$ exists and is unique up to natural isomorphism.
\end{theorem}
\begin{proof}
  % TODO: complete the proof.
\end{proof}

% ── Algebraic Properties ──────────────────────────────────────────────────────
\section{Algebraic Properties of the Lifted Torsion}
\label{sec:torsion-algebra}

% TODO: explore the algebraic consequences of the lifting, including
%       commutation relations and spectral theory.

% ── Summary ───────────────────────────────────────────────────────────────────
\section{Summary}
\label{sec:torsion-summary}

% TODO: summarize the lifting construction and its significance for
%       interpretations discussed in Chapter~\ref{chap:interpretations}.
